\section{Brainstorming(tormenta de ideas)}

\subsection{Objetivos y principios}

Brainstorming resuelve el \textbf{What \& Why}: «qué vamos a construir y por qué lo construimos así».

\begin{itemize}
    \item Se hace una sola pregunta de confirmación a la vez.
    \item Se entiende la intención real del usuario (Purpose), sus restricciones y los criterios de éxito.
    \item Primero se explora el contexto y se proponen \textbf{2--3 opciones}, analizando el Trade-off de cada una.
    \item Una vez que el usuario confirma, se forma y se guarda el \textbf{Design Doc} (\texttt{docs/plans/<date>-<topic>-design.md}).
\end{itemize}

\begin{warningbox}{Prohibiciones en la etapa de Brainstorming}
En la etapa de Brainstorming, \textbf{está prohibido invocar cualquier herramienta de implementación y prohibido escribir código}. El producto de esta etapa es puramente conceptual y de documentos de decisión, centrado en la arquitectura general, el flujo de datos, la interacción entre componentes y el manejo de excepciones.
\end{warningbox}

\subsection{Pasos del proceso}

\begin{enumerate}
    \item Explorar el contexto del proyecto.
    \item Plantear al usuario preguntas de confirmación para delimitar el problema con claridad.
    \item Proponer 2--3 soluciones y analizar el Trade-off.
    \item Presentar el Design Section y esperar la aprobación del usuario.
    \item Una vez aprobado por el usuario, invocar automáticamente el Skill \texttt{Writing Plans}.
\end{enumerate}

\subsection{Resumen de la sección}

El Brainstorming es la etapa obligatoria de esclarecimiento de requisitos; produce el Design Doc, que sienta la base arquitectónica para el plan de implementación posterior. Su disciplina central es «solo diseñar, no programar».


